\section{Espace $\Rn$}
\subsection{Espace vectoriel réel}
On désigne par $\Rn$ l'ensemble constitué de tous les $n$-tuples ordonnées $(x_1, \ldots, x_n)$ de nombres réels. Par la suite, les éléments de $\Rn$ seront notés indifféremment $\bm x$ ou $(x_1, \ldots, x_n)$. On munit alors $\Rn$ de deux opérations,
\[
 \forall \bm x, \bm y \in \Rn, \quad  \bm x + \bm y = (x_1 + y_1, \ldots, x_n + x_n)
\]
et 
\[
\forall \lambda \in \R, \forall \bm x \in \Rn, \quad \lambda \bm x = (\lambda x_1, \ldots, \lambda x_n)
\]
Avec ces deux opérations on vérifie immédiatement de $\Rn$ est un \emph{espace vectoriel} sur $\R$ de dimension $n$.

\subsection{Norme et produit scalaire}
Une des notions les plus importantes sur $\Rn$ est sans aucuns doute la notion de \emph{norme}

\begin{definition}[Norme] \label{def: norme}
	Une \emph{norme} est une application $N : \Rn \to \R_+$ telle que,
	\begin{itemize}
		\item $\forall \bm x \in \Rn, \quad N(\bm x) \geqslant 0 \et N(\bm x) = 0 \iff \bm x = \bm 0$
		\item $\forall \bm x \in \Rn, \forall \lambda \in \R, \quad N(\lambda \bm x) = |\lambda| N(\bm x)$
		\item $\forall \bm x, \bm y \in \Rn, \quad N(\bm x + \bm y) \leqslant N(\bm x) + N(\bm y)$
	\end{itemize}
	Plus communément on note une norme $\| \cdot \|$.
\end{definition}

\begin{definition}[Norme sur $\Rn$] \label{def: norme sur rn}
	On définit les application suivantes pour $p \in \R_+$, 
	\[
	\forall \bm x \in \Rn, \quad \norme{\bm x}_p = \left(\sum_{i=1}^n |x_i|^p \right)^{\frac{1}{p}}
	\]
	et
	\[
	\forall \bm x \in \Rn, \quad \norme{\bm x}_\infty = \max_{i \in \llbracket 1,n \rrbracket} |x_i|
	\]
\end{definition}

\begin{lemme}[Inégalité de Young] \label{lem: Inégalité de Young}
	\[
	\forall a,b \in \R_+, \quad a b \leqslant \frac{1}{p} a^p + \frac{1}{q} b^q
	\]
	où $p \in ]1, \infty[$ et $q \in \R$ tel que $\frac{1}{p} + \frac{1}{q} = 1$.
\end{lemme}

\prf{Soit $p \in ]1, \infty[$ et $q \in R$ tel que $\frac{1}{p} + \frac{1}{q} = 1$, par concavité de la fonction $t \mapsto \ln(t)$, on obtient
	\[
	\forall a,b \in \R_+, \quad \ln(ab) = \ln(a) + \ln(b) = \frac{1}{p} \ln(a^p) + \frac{1}{q} \ln(b^q) \leqslant \ln\left(\frac{1}{p} a^p + \frac{1}{q} b^q \right)
	\]
	ainsi par la croissance de la fonction $t \to \ln(t)$ sur $\R_+$ sur $\R_+$, on a bien,
	\[
	\forall a,b \in \R_+, \quad ab \leqslant \frac{1}{p}a^p + \frac{1}{q} b^q
	\]
}

\begin{lemme}[Inégalité de Hölder] \label{lem: Inégalité de Hölder}
	Soient $p \in ]1, \infty[$ et $q \in \R$ tel que $\frac{1}{p} + \frac{1}{q} = 1$, alors
	\[
	|\langle \bm x, \bm y \rangle | \leqslant \norme{\bm x }_p \norme{\bm y}_q
	\]
	où $\langle \cdot, \cdot \rangle$ désigne le produit scalaire usuel.
\end{lemme}

\prf{Si on pose,
	\[
	\forall i \in \entiers, \quad a_i = \frac{|x_i|}{\norme{\bm x}_p} \et b_i = \frac{|y_i|}{\norme{\bm y}_q}
	\]
	alors on obtient, par le lemme précédent \eqref{lem: Inégalité de Young},
	\[
	\sum_{i=1}^n a_i b_i \leqslant \frac{1}{p} \sum_{i=1}^n a_i^p + \frac{1}{q} \sum_{i=1}^n b_i^q = \frac{1}{p} + \frac{1}{q} = 1
	\]
	ainsi on obtient l'inégalité souhaitée,
	\[
	\sum_{i=1}^n \frac{|x_i||y_i|}{\norme{\bm x}_p \norme{\bm y}_q} \leqslant 1 \implies |\langle \bm x, \bm y \rangle | \leqslant \sum_{i=1}^n |x_i||y_i| \leqslant \norme{\bm x}_p \norme{\bm y}_q
	\]
}

\lecture{23 février}

\begin{theorem} 
	Les applications \eqref{def: norme sur rn} pour $p \in [1, \infty]$ définissent des normes sur $\Rn$.
\end{theorem}

\prf{Il faut distinguer trois cas,
	\begin{itemize}
		\item Supposons $p = \infty$, on a donc,
		\[
		\forall \bm x, \bm y \in \Rn, \quad \norme{\bm x + \bm y}_\infty = \max_{i \in \llbracket 1,n \rrbracket} |x_i + y_i| \leqslant \max_{i \in \llbracket 1,n \rrbracket} |x_i| + \max_{i \in \llbracket 1,n \rrbracket} |y_i| \leqslant \norme{\bm x}_\infty + \norme{\bm y}_\infty
		\]
		Et les autres propriétés \eqref{def: norme} d'une norme sont facilement vérifiées. 
		\item Supposons $p \in [1, \infty[$, pour $\bm x, \bm y \in \Rn$ on a,
		\[
		\norme{\bm x + \bm y}_p^p \leqslant \sum_{i=1}^n |x_i||x_i + y_i|^{p-1} + \sum_{i=1}^n |y_i||x_i + y_i|^{p-1}
		\]
		si $p=1$ on obtient directement l'inégalité triangulaire. Supposons donc $p \in ]1, \infty[$, on applique donc l'inégalité de Hölder \eqref{lem: Inégalité de Hölder} à l'inégalité précédente,
		\begin{align*}
			\norme{\bm x + \bm y}_p^p 	&\leqslant \sum_{i=1}^n |x_i||x_i + y_i|^{p-1} + \sum_{i=1}^n |y_i||x_i + y_i|^{p-1} \\
										&\leqslant \left( \sum_{i=1}^n |x_i|^p \right)^\frac{1}{p} \left(\sum_{i=1}^n |x_i + y_i|^q \right)^\frac{1}{q} + \left( \sum_{i=1}^n |y_i|^p \right)^\frac{1}{p} \left(\sum_{i=1}^n |x_i + y_i|^q \right)^\frac{1}{q} \\
										&\leqslant \norme{\bm x}_p \norme{\bm x + \bm y}_p^{p-1} + \norme{\bm y}_p \norme{\bm x + \bm y}_p^{p-1} \\
										&\leqslant (\norme{\bm x}_p + \norme{\bm y}_p) \norme{\bm x + \bm y}_p^{p-1}
		\end{align*}
		Si $\norme{\bm x + \bm y}_p = 0$ on a l'inégalité triangulaire, supposons donc $\norme{\bm x + \bm y}_p \neq 0$, ainsi en divisant de part et autre par $\norme{\bm x + \bm y}_p^{p-1}$, on obtient bien l'inégalité triangulaire souhaitée,
		\[
		\norme{\bm x + \bm y}_p \leqslant \norme{\bm x}_p + \norme{\bm y}_p
		\]
		Et les autres propriétés \eqref{def: norme} d'une norme sont facilement vérifiées.
		\item Finalement si nous supposons $p \in ]0,1[$. On pose,
		\[
		\bm x = (1, 0, \ldots, 0) \et \bm y = (0, \ldots, 0, 1)
		\]
		en supposant que l'inégalité triangulaire soit vraie, obtient
		\[
		2^\frac{1}{p} < 2
		\]
		ce qui est absurde car nous avions supposé que $p \in ]0,1[$.
	\end{itemize}
}

\begin{theorem}
	Toutes les normes $\{ \norme{\cdot}_p, ~p \in [1, \infty] \}$ sont équivalentes.
\end{theorem}

\prf{Soit $\bm x \in \Rn$, on observe que
	\begin{itemize}
		\item En posant $\bm y = (1,1, \ldots, 1) \in \Rn$ et en notant $|\bm x | = (|x_1, \ldots, |x_n|)$, on obtient
		\[
		\norme{\bm x}_1 = \sum_{i=1}^n |x_i| = \langle |\bm x|, \bm y \rangle \leqslant \norme{\bm x}_p \norme{\bm y}_q \leqslant n^\frac{1}{q} \norme{\bm x}_p
		\]
		\item Pour $p \in [1, \infty[$, on a
		\[
		\norme{\bm x }_p = \left(\sum_{i=1}^n |x_i|^p \right)^\frac{1}{p} \leqslant \left( n \left( \max_{i \in \llbracket 1,n \rrbracket} |x_i| \right)^p \right)^\frac{1}{p} = n^\frac{1}{p} \norme{\bm x}_\infty
		\]
		\item Et finalement,
		\[
		\norme{\bm x}_\infty = \max_{i \in \llbracket 1,n \rrbracket} |x_i| \leqslant \sum_{i=1}'n |x_i| = \norme{\bm x}_1
		\]
	\end{itemize}
	Ainsi on en déduit que n'importe quelle norme $\norme{\cdot}_\alpha$ pour être bornée par n'importe quelle autre norme $\norme{\cdot}_\beta$ multipliée par une constante $C_{\alpha, \beta} > 0$ indépendante de $\bm x$. Donc on a bien que toutes les normes $\{ \norme{\cdot}_p, ~p \in [1, \infty] \}$ sont équivalentes.
}



\begin{definition}
	Un \emph{produit scalaire} sur $\Rn$ est une application $b : \Rn \times \Rn \to \R$ telle que 
	\begin{enumerate} \label{def : produit scalaire}
		\item $\forall \bm x, \bm y \in \Rn, \quad b(\bm x, \bm y) = b (\bm y, \bm x)$
		\item $\forall \bm x, \bm y, \bm z, ~\alpha, \beta \in \R, \quad b(\alpha \bm x + \bm y, \bm z) = \alpha b(\bm x, \bm z) + b(\bm y, \bm z)$
		\item $\forall \bm x \in \Rn, \quad b(\bm x, \bm x) \geqslant 0$ et $b(\bm x, \bm x) = 0$ ssi $\bm x = 0$.
	\end{enumerate}
\end{definition}

\begin{proposition}
	On définit le \emph{produit scalaire usuelle} comme 
	\[
	\forall \bm x, \bm y \in \Rn, \quad \langle \bm x, \bm y \rangle = \bm x^\top \bm y = \bm y^\top \bm x = \sum_{i=1}^n x_i y_i
	\]
\end{proposition}

\prf{Il est facile de voir que cette application respecte bien tout les conditions d'un produit scalaire.}

\begin{lemme}
	Soit $b : \Rn \times \Rn$ un produit scalaire alors pour tout $\bm x, \bm y \in \Rn$ on a
	\[
	|b(\bm x, \bm y)| \leqslant \sqrt{b(\bm x, \bm x)}\sqrt{b(\bm y, \bm y)} = \norme{\bm x}_E \norme{\bm y}_E
	\]
	où $\norme{\cdot}_E$ désigne la norme induite par le produit scalaire $b$.
\end{lemme}

\prf{Pour tout $\bm x, \bm y \in \Rn$ et $\alpha \in \R$ on a
	\[
	0 \leqslant b(\alpha \bm x + \bm y, \alpha \bm x + \bm y) = \alpha^2 b(\bm x, \bm x) + 2\alpha b(\bm x, \bm y) + b(\bm y, \bm y) = p(\alpha)
	\]
	on a alors on polynôme positif ou nul en $\alpha$, on a donc un discriminant nécessairement négatif ou nul, on peut donc en conclure que
	\[
	4 b(\bm x, \bm y)^2 - 4 b(\bm x, \bm x) b(\bm y, \bm y) \leqslant 0
	\]
	D'où le lemme.
}
	

\section{Suites dans $\Rn$}

\begin{definition}[Suite convergente] \label{def: suite convergente}
	Soit $(\bm x^k) \subset \Rn$ une suite d'éléments de $\Rn$, $\bm x^k = (x_1^k, \ldots, x_n^k) \in \Rn$. On dit que $(\bm x^k)$ \emph{converge} s'il existe $\bm x \in \Rn$ tel que,
\[
\forall \varepsilon > 0, \exists N > 0, \forall k \geqslant N, \quad \norme{\bm x^k - \bm x} \leqslant  \varepsilon
\]
Dans ce cas on note $\lim_{k \to \infty} \bm x^k = \bm x$.
\end{definition}


\begin{proposition}[Indépendance de la norme] \label{prop: indépendance de la norme}
	Soient $(x^k) \subset \Rn $ une suite et $\norme{\cdot}_1, \norme{\cdot}_2$ deux normes sur $\Rn$, s'il existe $\bm x \in \Rn$ tel que,
	\[
	\lim_{k \to \infty} \norme{\bm x^k - \bm x}_1 = 0
	\]
	alors,
	\[
	\lim_{k \to \infty} \norme{\bm x^k - \bm x}_2 = 0
	\]
On a donc une indépendance de norme dans la caractérisation de la convergence d'une suite.
\end{proposition}


\prf{Puisque $\norme{\cdot}_1, \norme{\cdot}_2$ sont équivalente sur $\Rn$, on a
	\[
	\exists C > 0, \forall \bm x \in \Rn, \quad \norme{\bm x}_2 \leqslant C \norme{\bm x}_1
	\]
	Ainsi par hypothèse de convergence, pour $\varepsilon > 0$
	\[
	\exists N > 0, \forall k \geqslant N, \quad \norme{\bm x^k - x}_1 \leqslant \varepsilon
	\]
	Il suffit de poser $\epsilon =  \frac{1}{C} \varepsilon$, ainsi
	\[
	\forall k \geqslant N, \quad \norme{\bm x^k - \bm x}_2 \leqslant C \norme{\bm x^k- \bm x}_1 \leqslant C \varepsilon \leqslant \epsilon
	\]
}

\begin{lemme}[Convergence des suites composantes] \label{lem: convergence des suites composantes}
	Une suite $(\bm x^k) \subset \Rn$ converge vers $\bm x \in \Rn$ si et seulement si pour toute composante $j \in \entiers$ la suite $(x_j^k) \subset \R$ converge vers $x_j \in \R$.
\end{lemme}

\prf{Puisque la convergence ne dépend pas de la norme choisie \eqref{prop: Indépendance de la norme}, on considère en particulier la norme $\norme{\cdot}_\infty$. Par définition de la convergence d'une suite $(\bm x^k) \subset \Rn$
\[
\forall \varepsilon > 0, \exists N > 0, \forall k \geqslant N, \quad \norme{\bm x^k - \bm x}_\infty = \max_{i \in \llbracket 1,n \rrbracket} |x_i^k - x_i| \leqslant \varepsilon
\]
ce qui implique que $\lim_{k \to \infty} x_j^k \to x_j$ pour tout $j = 1, \ldots, n$.
Réciproquement, supposons qu'il existe $\bm x = (x_1, \ldots, x_n) \in \Rn$ tel que $\lim_{k \to \infty} x_j^k = x_j$ pour tout $j =1, \ldots, n$. Alors pour tout $j \in \entiers$,
\[
\forall \varepsilon > 0, \exists N_j, \forall k \geqslant N_j, \quad |x_j^k - x_j| \leqslant \varepsilon
\]
Si on pose $\tilde N = \max\{N_1, \ldots, N_n\}$, on a $\max_{i \in \llbracket 1,n \rrbracket} |x_i^k - x_i| \leqslant \varepsilon$ pour tout $k \geqslant \tilde N$, ce qui implique $\lim_{k \to \infty} \bm x^k = \bm x$.
}


\begin{definition}
	Une suite $(\bm x^k) \subset \Rn$ est dite de \emph{Cauchy} si
	\[
	\forall \varepsilon > 0, ~\exists N > 0, ~\forall k, l \geqslant N, \quad \norme{\bm x^k - \bm x^l} \leqslant \varepsilon
	\]	
\end{definition}


En suivant la même démonstration que le lemme \eqref{lem: convergence des suites composantes}, on peut montrer qu'une suite $(\bm x^k) \subset \Rn$ est de Cauchy si et seulement si chaque suite $(x_j^k) \subset \R$ pour $j = 1, \ldots, n$ est de Cauchy.


\begin{proposition}
	Une suite $(\bm x^k) \subset \Rn$ est convergente si et seulement si elle est de Cauchy.
\end{proposition}


\prf{On sait que cela est vrai en dimension $n = 1$. Ainsi,
	\begin{align*}
	(\bm x^k) \text{ converge} &\iff \forall j \in \entiers, \quad (x_j^k) \text{ converge} \\ &\iff \forall j \in \entiers, \quad (x_j^k) \text{ est de Cauchy} \iff (\bm x^k) \text{ est de Cauchy}
	\end{align*}
}

\begin{theorem}[Bolzano-Weierstrass] \label{thm: Bolzano-Weierstrass}
	Soit $(\bm x^k) \subset \Rn$ une suite \emph{bornée}. Alors il existe une sous-suite $(\bm x^{k_i}) \subset (\bm x^k)$ qui est convergente.
\end{theorem}

\prf{Puisque ${\bm x^k}$ est bornée, en particulier $(x_1^k)$ l'est aussi. Donc on peut extraire une sous-suite $(x_1^{k_j})$ qui converge vers $x_1 \in \R$. Désormais, comme la suite $(x_2^{k_j})$ est bornée, on peut extraire une sous-suite $(x_2^{k_{j_l}})$ qui converge vers $x_2 \in \R$. Ainsi,
	\[
	x_1^{k_{j_l}} \underset{l \to \infty}{\longrightarrow} x_1 \et x_2^{k_{j_l}} \underset{l \to \infty}{\longrightarrow} x_2
	\]

En itérant ce processus $n$ fois, on peut extraire une sous-suite de $(\bm x^k)$ donc chaque composante converge vers $(x_1, x_2, \ldots, x_n)$.
}

\section{Topologie de $\Rn$}

On s'intéresse ici à l'étude et classification de sous-ensembles de $\Rn$.
\begin{definition}[Boule] \label{def: Boule}
	Pour tout $\bm x \in \Rn$ et $\delta > 0$, on appelle
	\[
	B(\bm x, \delta) = \{ \bm y \in \Rn, \quad \norme{\bm x - \bm y} < \delta\}
	\]
	la \emph{boule ouverte} centrée en $\bm x$ et de rayon $\delta$. Ainsi que
	\[
	\overline B (\bm x, \delta) = \{ \bm y \in \Rn, \quad \norme{ \bm x - \bm y} \leqslant \delta \}
	\]
	la \emph{boule fermée} centrée en $\bm x$ et de rayon $\delta$.
\end{definition}
\lecture{25 février}
\begin{definition}[Ensembles ouverts/fermés] \label{def: ensembles ouverts/fermés}
	Soit $E \subset \Rn$. On dit que 
	\begin{itemize}
		\item $E$ est \emph{ouvert} si,
		\[
		\forall \bm x \in E, \exists \delta > 0, \quad B(\bm x, \delta) \subset E
		\]
		En particulier $E$ est ouvert si $E$ est vide.
		\item $E$ est \emph{fermé} si son complémentaire $E^c = \Rn\backslash E$ est ouvert.
	\end{itemize}
\end{definition}

\begin{proposition}
	Pour $\bm x \in \Rn$ et $\delta > 0$, $B(\bm x, \delta)$ est ouvert et $\overline B(\bm x, \delta)$ est fermée.
\end{proposition}

\prf{En effet, pour $\bm z \in B(\bm x, \delta)$,
	\[
	B(\bm z, \delta - \norme{\bm z - \bm x}) \subset B(\bm x, \delta)
	\]
	Donc $B(\bm x, \delta)$ est ouvert. De même, pour $\bm z \in \overline B(\bm x, \delta)^c$,
	\[
	B(\bm z, \norme{\bm z - \bm x} - \delta) \subset \overline B(\bm x, \delta)^c
	\]
	donc $\overline B(\bm x, \delta)$ est fermée.
}
\subsection{Classification des points de $\Rn$ par rapport à $E$.}
\begin{definition}
	Soit $E \subset \Rn$ et $\bm x \in \Rn$. On dit que
	\begin{itemize}
		\item $\bm x$ est un \emph{point intérieur} de $E$ si
		\[
		\exists \delta > 0, B(\bm x, \delta ) \subset E
		\]
		L'ensemble des points intérieurs de $E$ est noté $\mathring E$ et appelé \emph{intérieur} de $E$.
		\item $\bm  x$ est un \emph{point frontière} si,
		\[
		\forall \delta > 0, \quad B(\bm x, \delta) \cap E \neq \emptyset \et B(\bm x, \delta) \cap E^c \neq \emptyset
		\]
		l'ensemble des points frontière de $E$ est noté $\partial E$ et appelé la \emph{frontière} de $E$.
		\item $\bm x$ est un \emph{point adhérent} à $E$ si,
		\[
		\forall \delta > 0, \quad  B(\bm x, \delta) \cap E \neq \emptyset
		\]
		Un point adhérent est soit un point intérieur, soit un point frontière. L'ensemble des points adhérents à $E$ est noté $\overline E$, et appelé \emph{la fermeture} de $E$.
		\item $\bm x$ est un \emph{point isolé} de $E$ si,
		\[
		\exists \delta > 0, \quad B(\bm x, \delta) \neq E = \{ \bm x\}
		\]
		\item $\bm x$ est un \emph{point d'accumulation} si
		\[
		\forall \delta > 0, \quad B(\bm x, \delta) \cap (E \backslash \{\bm x\}) \neq \emptyset
		\]
		Il s'ensuit que les points d'accumulation de $E$ sont tous les points de $\overline E$ qui ne sont pas isolés.
	\end{itemize}
\end{definition}
\subsection{Quelques remarques sur les ensembles ouverts}
\begin{proposition}
	$\mathring E$ est ouvert. De plus $E$ est ouvert si et seulement si $E = \mathring E$.
\end{proposition}

\prf{Si $\bm x \in \mathring E$, il existe $\delta > 0$ tel que $B(\bm x, \delta) \subset E$. Or pour tout $\bm z \in B(\bm x, \delta)$, 
	\[
	B(\bm z, \delta - \norme{\bm z - \bm x}) \subset B(\bm x, \delta) \subset E
	\]
	et donc $\bm z \in \mathring E$ et $B(\bm x, \delta) \subset E$. \\
	On sait que $\mathring E \subset E$, or comme $E$ est ouvert on sait que pour tout $\bm x \in E$,
	\[
	\exists \delta > 0, \quad B(\bm x, \delta) \subset E
	\]
	donc $\bm x$ est un point intérieur de $E$ et donc $\bm x \in \mathring E$, ainsi $E = \mathring E$. La réciproque est trivial car $\mathring E$ est ouvert.
}

\begin{proposition}
	Toute réunion quelconque de sous-ensemble ouverts de $\Rn$ est un sous-ensemble ouvert. Cependant intersection finie de sous-ensembles ouverts de $\Rn$ est ouvert.
\end{proposition}

\prf{Soit $E = \bigcup_\alpha E_\alpha$ avec $E_\alpha$ ouvert. Pour tout $\bm x \in E$, il existe $\alpha$ tel que $\bm x \in E_\alpha$. Mais comme $E_\alpha$ est ouvert, alors il existe $\delta >  0$ tel que $B(\bm x, \delta) \subset E_\alpha \subset E$. Donc $E$ est ouvert. Similairement, soit $E = \bigcap_{i=1}^m E_i$ avec $E_i$ ouvert. Si $\bm x \in E$, alors pour tout $i =1,\ldots, m$ $\bm x \in E_i$, mais puisque chaque $E_i$ est ouvert, il existe $\delta_i > 0, \quad B(\bm x, \delta_i) \subset E_i$. On pose donc $\delta = \min \{\delta_i\}$ alors
	\[
	B(\bm x, \delta) \subset B(\bm x, \delta_i) \subset E, \quad \forall i
	\]
donc $B(\bm x, \delta) \subset \bigcap_{i=1}^m E_i = E$.
}

\subsection{Quelques remarques sur les ensembles fermés}

\begin{proposition}
	\[
	(\overline E)^c = \mathring{E^c} \et \overline{E^c} = (\mathring E)^c
	\]
\end{proposition}

\prf{On a
	\begin{align*}
		(\overline E)^c &= \{ \bm x \in \Rn ~|~ \exists \delta > 0, ~ B(\bm x, \delta) \cap E \neq \emptyset \} \\
		&= \{ \bm x \in \Rn ~|~ \exists \delta > 0, ~ B(\bm x, \delta) \subset E^c\} \\
		&= \mathring{E^c}
	\end{align*}
	Similairemment,
	\begin{align*}
		(\mathring E)^c &= \{\bm x \in \Rn ~|~ \forall \delta > 0, ~B(\bm x, \delta) \not \subset E\} \\
		&= \{\bm x \in \Rn ~|~ \forall \delta > 0, ~B(\bm x, \delta) \cap E^c \neq \emptyset \}\\
		&= \overline{E^c}
	\end{align*}
}

\begin{proposition}
	$\overline E$ est fermé. De plus $E$ est fermé si et seulement si $E = \overline E$.
\end{proposition}

\prf{Par passage on complémentaire on a
	\[
	(\overline E)^c = \mathring{E^c}
	\]
	qui est ouvert, donc $\overline E$ est fermé. De même
	\[
	E = \overline E \iff E^c = (\overline E)^c = \mathring E^c
	\]
	Donc $E^c$ est ouvert, donc $E$ est fermé.
}

\begin{proposition}
	Toute intersection quelconque de sous-ensembles fermé de $\Rn$ est un sous-ensemble fermé. Cependant toute intersection union finie de sous-ensembles fermés est fermée.
\end{proposition}

\prf{On sait que
	\[
	\left(\bigcap E_i\right)^c = \bigcup E_i^c \et \left( \bigcup E_i \right)^c = \bigcup E_i^c
	\]
	On conclut par le proposition similaire pour les sous-ensembles ouverts puisque si $E$ est fermée, alors $E^c$ est ouvert.
}

\begin{lemme}
	Un ensemble $E \subset \Rn$ non vide est fermé si et seulement si toute suite convergente $(\bm x^k) \subset E$ convergente, converge vers un élément de $E$.
\end{lemme}

\prf{Soit $E$ fermé et $(\bm x^k) \subset E$ une suite convergente vers $\bm x \in \Rn$. Alors $\bm x$ adhère $E$ et, puisque $E$ est fermé, $\bm x \in E$. Réciproquement, supposons que $E$ ne soit pas fermé, par passage au complémentaire $E^c$ n'est pas ouvert et donc il existe $\bm{\tilde x} \in E^c$ tel que $\forall \delta > 0, \quad B(\bm{\tilde x}, \delta) \not \subset E$. En posant $\delta = \frac{1}{k}$ pour tout $k \in \N^*$, on obtient alors une suite $(\bm x^k) \subset E$ convergente vers $\bm{\tilde x} \not \in E$.
}

\subsection{Ensembles compacts}

\begin{definition}
	Un ensemble $E \subset \Rn$ est \emph{compact} si et seulement s'il est à la fois borné et fermé.
\end{definition}
\lecture{3 mars}
\begin{theorem}[Caractérisation d'un sous-ensemble compact] \label{thm: Caractérisation d'un sous-ensemble compact}
	Un sous ensemble non vide $E \subset \Rn$ est compact si et seulement si de toute suite d'éléments de $E$ on peut extraire un sous-suite convergente qui converge vers un élément de $E$.
\end{theorem}


\prf{Soit $E$ un compact, alors par Bolzano-Weierstrass, de toute suite $(\bm x^k) \subset E$, on peut extraire un sous-suite $(\bm x^{k_j}) \subset E$ convergente vers $\bm x \in \Rn$ mais comme $E$ est fermé, alors $\bm x \in E$. \\
	Réciproquement, supposons que $E$ ne soit pas compact, donc si $E$ n'est pas fermé, il existe alors une suite $(\bm y^k) \subset E$ qui converge vers $\bm y \in E^c$. Or tout sous-suite extraite converge aussi vers $\bm y$. Si $E$ n'est pas borné, alors il existe un suite $(\bm x^k) \subset E$ telle que $\norme{\bm x^k} \geqslant k$ pour tout $k \in \N$. Toute sous-suite $(\bm x^{k_j})$ satisfait $\norme{\bm x^{k_j}} \geqslant k_j \geqslant j$ et donc la suite ne converge pas.
}

\subsection{Ensembles connexes et connexes pas arcs}

\begin{definition}
	Soit $E \subset \Rn$. On dit que $E$ est \emph{connexe} s'il existe pas deux ouvert $A,B \subset \Rn$ disjoint tels que $A \cap E \neq \emptyset, B \cap E \neq \et E \subset A \cap B$. Un ensemble non vide $E \subset \Rn$ est \emph{connexe par arcs} si pour tout $\bm x, \bm y \in E$, il existe un chemin continue $\bm \gamma : [0,1] \to E$ tel que $\bm \gamma(0) = \bm x$ et $\bm \gamma(1) = \bm y$.
\end{definition}
